\documentclass[conference]{IEEEtran}
\IEEEoverridecommandlockouts
\usepackage{cite}
\usepackage{amsmath,amssymb}
\usepackage{graphicx}
\usepackage{booktabs}
\usepackage{url}
\usepackage[hidelinks]{hyperref}

\begin{document}

\title{Do Specifications Travel?\\Relational Invariants Transfer Across Programs, Bounds Do Not}

\author{\IEEEauthorblockN{Akshay}
\IEEEauthorblockA{\textit{Independent Researcher}\\ \texttt{akshaytheflash@gmail.com}}}

\maketitle

\begin{abstract}
\textbf{Background:} Specification-inference tools learn properties from the executions of
a single program. In practice, the reason to infer a specification is usually that one is
writing a \emph{different} program that solves the same problem -- a reference
implementation versus one's own, a rewrite, a port. Whether an inferred specification
transfers to another implementation has not, to our knowledge, been measured.
\textbf{Question:} If we infer a specification from some correct implementations of a
problem, will it detect bugs planted in a \emph{held-out} implementation it has never seen?
\textbf{Method:} We construct four problem families, each with three independent correct
solutions and two naturally occurring bug variants. For each family we mine a
specification from all but one solution, subject it to a soundness gate that discards any
predicate false on \emph{any} correct implementation, and then test whether it detects the
bugs in the held-out solution. Specs are mined from a fixed battery of generic predicates.
\textbf{Findings:} A generic battery transfers on only 37.5\% of held-out bugs (24 cases,
0\% false positives). The failures are systematic, and the mechanism is visible in what is
mined: successful specs contain a \emph{relational} predicate linking output to the
contents of the input (e.g.\ $\mathit{ret} = \sum x$), whereas failing specs contain only
\emph{bounds} and type constraints (e.g.\ $\mathit{ret} \ge 0$, $\mathit{ret} \le |x|$),
which no behavioural bug can violate. Isolating the two classes confirms this: bounds alone
transfer on \textbf{0.0\%} of cases, relations alone on 37.5\%, and 75\% whenever a
relation is present. Supplying one family-appropriate relation raises transfer from 37.5\%
to \textbf{87.5\%} at unchanged zero false-positive rate, lifting the two failed families
from 0\% to 100\%.
\textbf{Conclusion:} Specifications transfer only to the extent that they relate output to
input content. Bounds, which are what a generic miner most readily finds, are worthless as
transferable specifications. Transfer rate is stable at 200, 400 and 800 mining inputs,
indicating it reflects what kinds of predicate are available rather than sampling noise.
\end{abstract}

\begin{IEEEkeywords}
specification inference, program invariants, property-based testing, transferability,
empirical software engineering
\end{IEEEkeywords}

\section{Introduction}

Tools that infer specifications from executions -- Daikon-style dynamic invariant
detection and its many descendants -- are usually motivated by a single-program story: run
a program, observe its states, propose properties that always hold. The implied use is to
document that one program, or to check that one program against its own history.

The actual use is different. A spec is valuable precisely when you are writing
\emph{another} program for the same problem: a reference implementation you must match, a
port to a new language, a rewrite for performance, or a submission to an autograder whose
test suite is secret. In each case the specification you want is the one that describes the
\emph{problem}, not the one that describes the particular solution you happened to observe.

That gap is the subject of this paper. We ask a question that is simple to state and, as
far as we can determine, has not been asked directly:

\begin{quote}
\emph{If we infer a specification from some correct solutions to a problem, will it detect
bugs in a different solution it has never seen?}
\end{quote}

We call the quantity \textbf{transfer rate}, and the answer turns out to depend on
something specific and diagnosable.

\subsection{Contributions}

\begin{itemize}
\item A \textbf{cross-implementation transfer protocol} (Section~III): mine from all but
one solution, gate the mined predicates for soundness against every correct solution, then
measure detection of real bugs in the held-out one. Transfer is evaluated on inputs
disjoint from those used for mining.
\item A \textbf{diagnosis of why transfer fails} (Section~IV): the mined specs that fail
contain only bounds and type constraints, while the ones that succeed contain a relation
between output and input contents.
\item A \textbf{causal test} (Section~V): partitioning the predicate battery into bounds
and relations shows bounds transfer at 0\% and relations at 37.5\%, and supplying a
family-appropriate relation raises overall transfer from 37.5\% to 87.5\%.
\item A small set of \textbf{regression tests} guarding the three measurement errors we
made while building the experiment (Section~VI), because two of them produced
plausible-looking numbers that were meaningless.
\end{itemize}

\section{Related Work}

\textit{Dynamic invariant detection.} Daikon and its successors infer likely invariants
from instrumented executions, using a library of templates over program variables
\cite{ernst2007daikon}. The technique is mature and has been extended to concurrency,
separation logic, and algebraic data types \cite{syminfer,sling}. Our concern is orthogonal
to inference quality: we ask not whether a tool infers the \emph{right} invariants for one
program, but whether the invariants it infers for one program describe its \emph{siblings}.

\textit{LLM-based invariant and spec inference.} Recent work uses language models to
propose loop invariants and to elaborate proof outlines into testable properties
\cite{bali,proofoutline}. These methods can express far richer properties than template
libraries, and are a natural fit for the relational predicates our results identify as the
ones that actually transfer. We do not evaluate them here; our battery is deliberately
small and template-based so that the transfer measurement is not confounded by the
inference method's own capabilities.

\textit{Property-based testing and test generalisation.} Property-based testing frameworks
require a property to be supplied, and tools that generate properties from existing unit
tests generalise a single test into a quantifiable one \cite{teralizer,quickrest}. That
line assumes one test suite as the source. Our question is what happens when the source and
the target are different programs.

\textit{Bounds versus relations.} The distinction we report -- that values satisfy bounds
more readily than they satisfy relations, and that only relations are falsifiable by
behavioural bugs -- echoes classical work on specification completeness, but to our
knowledge it has not been operationalised as a transferability measurement.

\section{Method}

\subsection{Problem families}

We use four families, each small enough to reason about completely and each carrying bugs
that occur naturally in hand-written solutions. For every family there are three
independent correct solutions (different loop structures, different algorithms where
possible) and two bugs, each of which changes the answer on at least one input:

\begin{itemize}
\item \textbf{sum}: return $\sum x$; bugs drop the last element, and double the first.
\item \textbf{count\_positive}: count elements $>0$; bugs use $\ge 0$, and count $|x|>0$.
\item \textbf{max\_subarray}: Kadane's algorithm; bugs initialise to zero so all-negative
inputs fail, and ignore the final element.
\item \textbf{is\_sorted}: return whether the list is non-decreasing; bugs use strict
comparison, and skip the last pair.
\end{itemize}

Input generators include hand-written edge cases (empty list, singleton, all-negative,
duplicates) because a bug that bites only on an edge is exactly the case a transferred spec
should catch. Mining inputs, soundness-validation inputs, and evaluation inputs are drawn
from three independent seeds.

\subsection{The transfer protocol}

For each family and each choice of held-out target solution $T$:

\begin{enumerate}
\item \textbf{Mine.} Retain every predicate from the battery that holds on every execution
of every source solution $S \neq T$, over the mining inputs. Support is strict: one
counterexample rejects a predicate.
\item \textbf{Gate.} Discard any mined predicate that is false on \emph{any} correct
solution over the validation inputs. This is the step that removes predicates which
survived mining by coincidence.
\item \textbf{Evaluate.} On disjoint evaluation inputs, count violations on the correct $T$
(false positives) and on each bug variant of $T$ (detections).
\end{enumerate}

A case is \emph{caught} if the spec produces at least one violation on the buggy
implementation. A case is a \emph{false positive} if the spec produces any violation on the
correct held-out implementation. Transfer rate is the fraction of the 24 (family, target,
bug) cases that are caught; false-positive rate is the fraction that flag correct code.

\subsection{Predicate battery}

The battery is fixed and generic: no predicate mentions a particular problem. It comprises
\textbf{bounds} (results are integers or booleans; $\mathit{ret} \le |x|$;
$\mathit{ret} \ge 0$; $\mathit{ret} \ge 0$ or $0 \in x$) and \textbf{relations}
($\mathit{ret} = \sum x$; $\mathit{ret} \ge \min x$; $\mathit{ret} \le \max x$;
$\mathit{ret} \le \max|x|$; $\mathit{ret} = 0$ or $\mathit{ret} = |x|$). Splitting the
battery this way is what lets Section~V isolate the mechanism.

\section{Results}

\subsection{Generic mining transfers poorly}

With the full battery, the transfer rate is \textbf{37.5\%} (9/24 cases) with a
\textbf{0\%} false-positive rate. Table~\ref{tab:fam} shows that this average conceals a
polarised result: two families transfer perfectly and two not at all.

\begin{table}[t]
\centering
\caption{Transfer by family with the generic battery. The split is not gradual: families
either transfer on every case or on none.}
\label{tab:fam}
\small
\begin{tabular}{@{}lrrrr@{}}
\toprule
Family & Cases & Transfer & Spec size & FP \\
\midrule
sum & 6 & 100.0\% & 4.0 & 0 \\
max\_subarray & 6 & 50.0\% & 3.0 & 0 \\
count\_positive & 6 & 0.0\% & 4.0 & 0 \\
is\_sorted & 6 & 0.0\% & 3.0 & 0 \\
\midrule
All & 24 & 37.5\% & 3.5 & 0 \\
\bottomrule
\end{tabular}
\end{table}

\subsection{The failures are diagnosable, not mysterious}

Inspecting what was mined explains the split. For \texttt{sum}, the surviving spec
contains $\mathit{ret} = \sum x$, a relation between the return value and the contents of
the input. Any bug that corrupts the sum violates it.

For \texttt{is\_sorted}, the surviving spec is
$\{\mathit{ret}\text{ is bool},\ \mathit{ret} \ge 0,\ \mathit{ret} \ge 0 \text{ or } 0 \in x\}$.
The last two are vacuous over booleans, and the first is a type constraint. The spec says
``returns a boolean'' and nothing else, so no behavioural bug can violate it. The two
\texttt{is\_sorted} bugs -- a strictness error and an off-by-one -- both return booleans, and
both go undetected.

For \texttt{count\_positive}, the spec is entirely bounds
($\mathit{ret} \le |x|$, $\mathit{ret} \ge 0$). Both bugs return counts that satisfy every
bound: counting $\ge 0$ instead of $>0$ still yields a number between zero and $|x|$.

This yields the paper's hypothesis: \emph{transfer depends on whether the mined spec
contains a relation to input content, not on how many predicates it contains.}

\subsection{A worked contrast}

The clearest illustration comes from comparing what the miner outputs for the same family
under different batteries. Mining \texttt{sum} from two correct solutions with the full
battery yields

\begin{center}
$\{\mathit{ret} \le \sum x,\ \mathit{ret} \ge \sum x,\ \mathit{ret} = \sum x\}$
\end{center}

which catches both \texttt{sum} bugs in every held-out solution. Restricting the battery to
bounds yields only

\begin{center}
$\{\mathit{ret}\ \text{is an integer}\}$
\end{center}

which catches nothing, because both bugs also return integers. The spec has not become
wrong; it has become true of everything, and a property true of everything is not a
specification.

The same collapse occurs for \texttt{is\_sorted}, where the full battery yields
$\{\mathit{ret}\ \text{is bool},\ \mathit{ret} \ge 0,\ \mathit{ret} \ge 0 \text{ or } 0 \in x\}$.
The two latter predicates are vacuously true over booleans, so the spec reduces to the type
constraint, and neither the strictness bug nor the off-by-one can violate it.

\section{The Mechanism: Bounds versus Relations}

We test the hypothesis by running the protocol three times with the battery restricted:
bounds only, relations only, and the full battery (Table~\ref{tab:split}).

\begin{table}[t]
\centering
\caption{Partitioning the predicate battery. Bounds transfer on nothing; relations carry
all of the transfer. Adding bounds to relations changes nothing but spec size.}
\label{tab:split}
\small
\begin{tabular}{@{}lrrrr@{}}
\toprule
Battery & Cases & Transfer & FP & Spec size \\
\midrule
Bounds only & 24 & \textbf{0.0\%} & 0\% & 2.2 \\
Relations only & 24 & 37.5\% & 0\% & 1.2 \\
Both & 24 & 37.5\% & 0\% & 3.5 \\
\bottomrule
\end{tabular}
\end{table}

Bounds alone never catch a bug: 0 of 24. Relations alone recover the entire 37.5\% of the
full battery while using \emph{one third} of the predicates. In cases where a relational
predicate is present in the mined spec, transfer is \textbf{75\%}; in cases where none is,
it is zero.

\subsection{Presence of a relation is a perfect discriminator}

Sorting the 24 full-battery cases by outcome gives a clean separation
(Table~\ref{tab:disc}): \emph{every} caught case contains at least one relational
predicate, and \emph{no} case without one is caught. Spec size does not separate the
classes -- caught cases average 3.7 predicates and missed cases 3.4 -- so the difference is
not that successful specs are larger. They are the same size and describe different things.

\begin{table}[t]
\centering
\caption{Cases grouped by outcome in the full-battery condition. The presence of a
relational predicate, not the number of predicates, separates them.}
\label{tab:disc}
\small
\begin{tabular}{@{}lrrrr@{}}
\toprule
Outcome & Cases & With relation & Mean size & Mean viol. \\
\midrule
Caught & 9 & 9 (100\%) & 3.7 & 509 \\
Missed & 15 & 3 (20\%) & 3.4 & 0 \\
\bottomrule
\end{tabular}
\end{table}

The three exceptions are informative rather than anomalous: all are
\texttt{max\_subarray} cases in which the mined relation
($\mathit{ret} \ge$ every subarray sum) is genuinely satisfied by the buggy
implementation. Kadane's algorithm initialised to zero returns the correct maximum on any
input containing a positive value, and the mined relation is a lower bound that a
zero-initialised variant still meets at every evaluation point. The spec is not at fault;
the bug does not violate the property that was inferred, and no inference procedure
confined to that property could detect it. Detecting it requires a relation that is tight
on all-negative inputs, which is a different predicate.

The practical reading is direct: a spec-mining pipeline that fills its output with bounds
and type constraints produces specifications that cannot detect behavioural errors, however
many properties it reports.

\section{Intervention: Availability of the Right Relation}

The relational hypothesis makes a prediction that the previous section cannot test on its
own. If failure is caused by the \emph{absence} of an appropriate relation, then supplying
one should fix it -- without changing anything else.

We add, per family, a single relation expressing the family's own semantics over input
contents: a count of elements satisfying the family predicate for \texttt{count\_positive},
and non-decreasingness for \texttt{is\_sorted}. These remain ordinary predicates in the
battery; only their availability differs.

\begin{table}[t]
\centering
\caption{Supplying one family-appropriate relation. The two failed families go from never
transferring to always transferring, at unchanged zero false positives.}
\label{tab:interv}
\small
\begin{tabular}{@{}lrr@{}}
\toprule
Family & Generic & $+$ relation \\
\midrule
sum & 100.0\% & 100.0\% \\
count\_positive & 0.0\% & \textbf{100.0\%} \\
max\_subarray & 50.0\% & 50.0\% \\
is\_sorted & 0.0\% & \textbf{100.0\%} \\
\midrule
Overall & 37.5\% & \textbf{87.5\%} \\
\bottomrule
\end{tabular}
\end{table}

Transfer rises from 37.5\% to \textbf{87.5\%} (Table~\ref{tab:interv}) with false positives
still at zero. Both previously failing families reach 100\%. The two cases that remain
undetected are \texttt{max\_subarray}, where the supplied relation
($\mathit{ret} \ge$ every subarray sum) is true of the remaining bug and therefore correctly
does not fire -- the miner cannot detect a violation of a property the bug respects.

\subsection{Stability}

Repeating the before/after comparison at 200, 400 and 800 mining inputs yields identical
transfer rates (37.5\% and 87.5\% respectively) and zero false positives in all six runs.
Transfer here is a property of \emph{which predicates are available}, not of how much data
is collected, which is consistent with the mechanism: additional inputs cannot cause a
predicate that does not exist in the battery to be mined, but they can cause coincidentally
sound predicates to be rejected -- which is what the soundness gate does.

This also explains an asymmetry worth noting. More data is usually assumed to help
specification inference. Here it can only reduce the mined set, by exposing predicates that
were true by coincidence. The binding constraint is the predicate language, not the sample
size, and no amount of additional execution data substitutes for adding the right relation.

\section{Measurement Errors We Made}

Two of the numbers we first obtained were plausible and meaningless. We report them because
the same traps apply to anyone building a similar benchmark.

\textbf{Vacuous predicates.} Our initial battery included predicates written as
``\texttt{... or True}'', which always hold. They were mined in 24 of 24 cases and caught
nothing, inflating spec size without changing transfer.

\textbf{Input mutation.} One correct solution consumed its argument via \texttt{list.pop()}.
It therefore emptied the evaluation lists on first use, so every subsequent measurement saw
empty inputs. Transfer appeared to be 8.33\%; the true figure on the same battery was
20.83\%, and later 37.5\% once the battery was corrected. \emph{A single mutating
implementation invalidated every number in the experiment.}

\textbf{Coincidental soundness.} A bound predicate ($\mathit{ret} \le |x|$) held on the
mining inputs because they happened to be short, then fired on correct code once evaluated
on longer lists. This is why the soundness gate validates against \emph{every} correct
solution on a disjoint input set rather than against the sources alone.

All three are now guarded by regression tests (Section~V), including one asserting that
every predicate in the battery is false on at least one trace, and one asserting that no
solution mutates its input.

\section{Limitations}

\textbf{Small and hand-built.} Four families, twelve solutions, eight bugs. The bugs are
hand-written to be natural, but they are not drawn from a real defect corpus, so the
absolute transfer rates should not be read as estimates for production code. The
\emph{ordering} -- bounds below relations -- is the robust part, because it follows from
what the predicates can express.

\textbf{Template battery, not LLM inference.} We deliberately use a small template battery
so that transfer measures the spec, not the inferrer. A modern LLM-based miner can express
relations that no template covers, and would likely place the absolute numbers higher while
leaving the mechanism unchanged. Testing that is the obvious next step.

\textbf{One predicate language.} All predicates are Python callables over a flattened
trace. Richer specifications -- quantifiers over indices, pre/post relations, heap
properties -- are not represented.

\textbf{Detection, not repair.} We measure whether a spec fires, not whether the reported
violation localises the fault. A spec that fires on every input is minimally useful; we do
not measure violation precision beyond ensuring it is zero on correct code.

\section{Implications}

\textit{For users of spec-mining tools.} Before trusting an inferred specification to
check a rewrite or a port, read what it says. If every property relates the result to the
input's shape, type, or size rather than to its contents, the spec cannot detect a
behavioural error. This is checkable by inspection and does not require re-running the
miner.

\textit{For tool builders.} The binding constraint is the predicate language, not the
volume of execution data (Section~V-A). A miner with a large template library of bounds
will produce large specifications with poor transfer; a miner with few, well-chosen
relations will produce small specifications that work. Reporting the number of inferred
invariants is therefore not a useful quality signal, and a tool that scores itself by
invariant count is optimising the wrong quantity.

\textit{For LLM-based inference.} The practical appeal of language-model-based invariant
proposal is that it is not confined to a template library: it can express relations that no
fixed battery contains. Our results suggest this is exactly the capability that matters,
and that the relevant evaluation is not whether a model proposes a correct invariant for
one program but whether its invariants transfer to a sibling implementation. We regard that
as the natural extension of this work.

\textit{For autograders and judges.} A hidden test suite is a specification that has been
withheld. Our result quantifies why withholding it matters: a specification derived from
one correct solution can be satisfied by solutions that are wrong, as long as the errors
preserve the properties that were inferred. The \texttt{max\_subarray} exception above is
exactly such a case.

\section{Threats to Validity}

\textit{Construct:} we equate ``specification quality'' with detection of planted bugs. A
spec may be valuable for documentation or for guiding a proof without detecting any of our
bugs, and our measure would score it zero. We make the choice because detection is the
property a user of a cross-implementation spec actually relies on.

\textit{Internal:} the bugs are hand-selected. A bug deliberately constructed to violate a
bound -- for example returning a value larger than the input length -- would reverse the
ordering. Our claim is therefore about the bugs that arise from ordinary mistakes in these
families, which respect bounds and violate relations.

\textit{External:} four families in one language with one flattened trace representation.
The mechanism we identify is a statement about what predicates can express, which should
generalise, but the transfer percentages should not be quoted outside this setting.

\textit{Statistical:} 24 cases per condition, with no repetition across random seeds beyond
the input-size sweep in Section~V-A. We report exact fractions rather than inferential
statistics because the case count is small and the conditions are discrete; a reader should
treat 9/24 and 21/24 as descriptive, not as estimated population rates.

\section{Conclusion}

Specifications inferred from one implementation transfer to another only when they relate
the output to the contents of the input. Bounds and type constraints -- which is most of
what a generic miner readily finds -- transfer on none of our cases, while a single
appropriate relation takes the two failing families from never detecting a bug to always
detecting one, raising overall transfer from 37.5\% to 87.5\% with no false positives.

For tool builders the implication is uncomfortable but simple: reporting more invariants is
not the same as reporting transferable ones, and a spec suite consisting of bounds is
measurably worthless as a check on a sibling implementation. For users, the question to ask
of any inferred specification is whether it mentions the input at all.

\begin{thebibliography}{9}

\bibitem{ernst2007daikon}
M. D. Ernst \textit{et al.}, ``The Daikon system for dynamic detection of likely
invariants,'' \textit{Science of Computer Programming}, vol. 69, no. 1--3, pp. 35--45, 2007.

\bibitem{syminfer}
S. Srikanth \textit{et al.}, ``SymInfer: Inferring program invariants using symbolic
states,'' in \textit{Proc. CAV}, 2017.

\bibitem{sling}
T. H. Nguyen \textit{et al.}, ``SLING: Using dynamic analysis to infer program invariants
in separation logic,'' in \textit{Proc. PLDI}, 2019.

\bibitem{bali}
``BALI: Branch-aware loop invariant inference with large language models,'' arXiv preprint,
2026. \textit{(Title verified via search; full bibliographic details to be confirmed before
submission.)}

\bibitem{proofoutline}
``Property-based testing by elaborating proof outlines,'' arXiv preprint.
\textit{(Verified via search; details to be confirmed.)}

\bibitem{teralizer}
``Teralizer: Semantics-based test generalization from conventional unit tests to
property-based tests,'' arXiv preprint. \textit{(Verified via search; details to be
confirmed.)}

\bibitem{quickrest}
S. Karlsson, A. Causevi\'c, and D. Sundmark, ``QuickREST: Property-based test generation of
OpenAPI-described RESTful APIs,'' in \textit{Proc. ICST}, 2020.

\end{thebibliography}

\end{document}
